\begin{abstract}
% SYNC-BEGIN ABS-01
Neuronal morphology analysis supports reconstruction proofreading, connection completion, and cell characterization from local neurites to whole cells. Jointly learning from skeletons and surfaces faces two linked obstacles: limited sampling can omit geometry and obscure size differences, while object-level fusion lacks an explicit association between a branch and its surrounding surface. We formulate these obstacles through measure-based geometric representation: sampling determines which geometry is represented and with what weight, and regional aggregation determines where heterogeneous evidence is combined. Guided by this formulation, we propose the Common-Region Skeleton--Surface (CRSS) framework with two mechanisms. Weighted multiresolution geometric encoding combines global and local surface samples with the observed skeleton and position/length embeddings. Topology-guided common-region fusion separately aggregates skeletal and surface evidence on shared arclength regions and learns joint tokens before contextual propagation and task readout. We evaluate the framework on five tasks from NREDD, H01, and FlyWire using matched development-set protocols with seed 42. At the same point budget, the sampling mechanism improves FlyWire macro-F1 by 9.91 percentage points over uniform sampling; learned regional interaction improves H01 local and FlyWire classification by 1.92 and 8.90 points over regional averaging. These results support connecting input coverage with explicit local correspondence for neuronal morphology learning.
% SYNC-END ABS-01
\end{abstract}
